% !TEX root = ../thesis.tex

\section{Synthetic Data Generation Methods}
\label{sec:method:generators}

% Expansion of proposal §5.2 (Proposed Models). One subsection per generator
% with architecture, hyperparameters, training protocol, and (where
% applicable) the differential-privacy mechanism.

This section specifies the five generators whose synthetic outputs are
compared throughout the thesis. The rationale for the selection is set out
first (\S\ref{sec:methods:rationale}), followed by a full specification of
each generator in turn: the HealthGAN WGAN-GP baseline
(\S\ref{sec:methods:healthgan}), the differentially private PATE-GAN
(\S\ref{sec:methods:pategan}), and the three language-model synthesisers built
on the Pythia family at two scales and under a private regime
(\S\S\ref{sec:methods:pythia}-\ref{sec:methods:pythia-dp}). Each generator is
described by its architecture, its encoding or serialisation scheme, the
training protocol, and, where applicable, the differential-privacy mechanism.

\subsection{Selection Rationale}
\label{sec:methods:rationale}

The five generators evaluated in this thesis populate a comparative design
space spanned by two axes. The first is the generative paradigm, which
contrasts GANs with LLMs. The second is the privacy regime, which contrasts models
trained without a formal guarantee against models trained under DP.
On the LLM side a third axis, model scale, is introduced via a small
(Pythia-70M) and a large (Pythia-1B) non-DP baseline. The DP variant is
applied at the larger scale, for reasons set out in
\S\ref{sec:methods:pythia-dp:why-1b}. This arrangement permits the effect
of formal privacy to be isolated within each paradigm, the two paradigms
to be compared under a common privacy regime, and the LLM-side scale
sensitivity to be characterised. The rationale for each generator is set
out below.

\paragraph{HealthGAN:}
HealthGAN couples the medGAN design with the Wasserstein GAN with gradient
penalty (WGAN-GP) framework and a dedicated transformation pipeline for the
mixed continuous and categorical attributes characteristic of electronic
health records. HealthGAN also introduces Nearest-Neighbour Adversarial Accuracy, which
evaluates both data resemblance and privacy. It achieves competitive
utility and resemblance compared with established baselines and supports
the secure export of trained models~\cite{healthgan}. These properties make it a
suitable representative of the non-private GAN paradigm
(\S\ref{sec:bg:gan:health}).

\paragraph{PATE-GAN:}
PATE-GAN provides formal $(\varepsilon,\delta)$-DP by adapting the Private
Aggregation of Teacher Ensembles framework~\cite{Papernot2017PATE} to
generative modelling. An ensemble of teacher discriminators is trained on
disjoint data partitions, their predictions are combined through a
differentially private noisy aggregation mechanism, and a student
discriminator is trained on the resulting labels. The generator
inherits the privacy guarantee by post-processing~\cite{pategan}. It thus
supplies a GAN-based generator whose privacy budget is explicitly controllable
and accountable (\S\ref{sec:bg:gan:pategan}).

\paragraph{Pythia-70M:}
The benchmark study evaluates the Pythia family of decoder-only language
models for synthetic tabular health data. It uses the
Tabula framework, where each table row is serialised as a text sequence.
The study reports that Pythia-based generators achieve utility that matches
or exceeds CTGAN~\cite{Miletic2024}. Pythia-70M is included here as the
small-scale non-private LLM representative (\S\ref{sec:bg:llm:tabula}), and,
as discussed in \S\ref{sec:methods:pythia-dp:why-1b}, also frames the
model-scale precondition that motivated migrating the LLM-side DP variant to
the 1B model.

\paragraph{Pythia-1B:}
A non-private Pythia-1B variant is included to provide a more closely
matched comparison with Pythia-1B-DP. Keeping the model size fixed avoids
combining the effect of DP-SGD with the increase from $70$M to $1$B
parameters, although other training differences remain. Pythia-1B
uses the same Tabula serialisation and the same LoRA fine-tuning protocol as
Pythia-70M, differing only in the underlying decoder-only backbone.

\paragraph{Pythia-1B-DP:}
Beyond the generators specified in the original proposal, this thesis
introduces a differentially private variant of Pythia-1B trained with
differentially private stochastic gradient descent
(DP-SGD)~\cite{Abadi2016DPSGD}, as implemented in the Opacus
library~\cite{Yousefpour2021Opacus}. This addition completes the design space
and enables a within-paradigm assessment of the cost of formal privacy for the
LLM, paralleling the HealthGAN-to-PATE-GAN comparison on the GAN side.

\subsection{HealthGAN: WGAN-GP Baseline}
\label{sec:methods:healthgan}

% Code references:
%   - SDV encoding/decoding: ../../healthgan/generators/sdv_converter.py
%   - WGAN-GP training:      ../../healthgan/generators/wgan_for_mac.py

HealthGAN realises the non-private GAN baseline as a Wasserstein GAN with
gradient penalty (WGAN-GP) trained on a fully numeric encoding of the
diabetes dataset~\cite{healthgan}
.\footnote{The implementation used in this
thesis was adapted from the original HealthGAN GitHub
repository~\cite{yale2019healthganCode},
\url{https://github.com/yknot/ESANN2019/tree/master}, and modified for the
dataset and experimental pipeline used in this study.} The raw mixed-type table is first mapped
to a continuous representation bounded in $[0,1]$. A generator and a critic
are then trained adversarially on that representation; and synthetic records
are recovered by inverting the encoding. The subsubsections below describe,
in turn, the encoding pipeline, the network architecture, the training
objective, the training configuration, and the selection of the synthetic
sample carried forward to the comparative evaluation.

\subsubsection{SDV Encoding and Decoding}
\label{sec:methods:healthgan:sdv}

The networks operate not on the raw table but on an encoded representation in
which every column is mapped to the unit interval $[0,1]$, yielding a purely
continuous and bounded input. The encoder dispatches each column to one of
four transformations according to its data type~\cite{patki2016synthetic}.

\paragraph{Categorical columns:}
Categorical variables are first ordered by decreasing relative frequency, with
category names used to break ties. Each category is then assigned a
sub-interval of $[0,1]$, where the width of the interval corresponds to that
category's relative frequency in the column. Each observed categorical value
is replaced by a random draw from a truncated normal distribution within its
assigned interval. The distribution is centered at the midpoint of the
interval, and its standard deviation is set to one sixth of the interval
width. For decoding, the cumulative interval boundaries and their
corresponding categories are stored as a lookup table.

\paragraph{Ordinal columns:}
Non-binary integer columns are treated as ordinal variables and encoded using
the same interval-partitioning scheme. Unlike nominal categorical variables,
the intervals are ordered by the integer values rather than by category
frequency. Any integer levels missing between the observed minimum and maximum
are added with zero count. Additive smoothing with $\alpha = 1$ is then
applied to the level frequencies before constructing the intervals, ensuring
that even unobserved levels receive a non-zero interval width.

\paragraph{Binary columns:}
Columns containing only 0 and 1 are treated as categorical variables by
default and are encoded using the same interval-based transformation. If the
optional \texttt{-beta} flag is enabled, they are instead encoded using
truncated Beta distributions whose parameters are derived from the observed
class balance.

\paragraph{Numeric columns:}
Floating-point columns are min-max scaled to $[0,1]$, with the original
minimum and maximum stored for decoding. Low-cardinality coded variables,
such as the binned age field, are forced onto the ordinal path so that they
are not treated as continuous quantities.

Missing values are resolved before transformation. Categorical columns
receive a placeholder level, ordinal and binary columns are filled with the
mode, and numeric columns with the median.

Decoding inverts each transformation. Interval-encoded columns are mapped
back to their corresponding categories, while min-max columns are rescaled
to their original ranges. Because the generator can emit values marginally
outside $[0,1]$, generated values are first clipped to that interval. A
final snap-to-nearest step then restores schema compliance. A numeric column
with at most fifty distinct observed values is snapped to its nearest valid
value, while a high-cardinality numeric column is clipped to a non-negative
integer.

\subsubsection{WGAN-GP Architecture}
\label{sec:methods:healthgan:arch}

Both networks are fully connected feed-forward networks whose layer widths
are defined relative to the encoded feature count $F$ and a base width $B$.
For the encoded diabetes table $F = 45$, and the configuration used in this
thesis sets $B = 64$.

The generator transforms a latent noise vector into a synthetic record. Its
input is a $100$-dimensional vector drawn from a standard multivariate normal
distribution, $z \sim \mathcal{N}(\mathbf{0}, \mathbf{I}_{100})$, which is
passed through three dense layers:
\begin{equation}
  100 \;\to\; 2F~(\mathrm{ReLU}) \;\to\;
  \mathrm{round}(1.5F)~(\mathrm{ReLU}) \;\to\; F~(\mathrm{sigmoid}),
\end{equation}
that is, hidden widths of $90$ and $68$ units and an output of $45$ units.
The terminal sigmoid constrains every output to $[0,1]$, matching the range
of the encoded representation.

The critic maps a record to an scalar score through four dense
layers:
\begin{equation}
  F \;\to\; B~(\mathrm{LeakyReLU}) \;\to\; 2B~(\mathrm{LeakyReLU})
  \;\to\; 4B~(\mathrm{LeakyReLU}) \;\to\; 1,
\end{equation}
with hidden widths of $64$, $128$, and $256$ units and a single linear output
unit. Consistent with the Wasserstein formulation, the output is left
unactivated. The critic estimates a real-valued distance rather than a
classification probability.

\subsubsection{Training Loss: Wasserstein with Gradient Penalty}
\label{sec:methods:healthgan:loss}

The networks are trained under the Wasserstein GAN objective with gradient
penalty~\cite{Arjovsky2017WGAN, Gulrajani2017WGANGP}. Let $G$ denote the
generator and $D$ the critic, and let $P_r$ and $P_g$ denote the real and
generated data distributions, respectively. The critic is trained to
minimise the Wasserstein objective with gradient penalty,
\begin{equation}
  \mathcal{L} =
    \underbrace{\mathbb{E}_{\tilde{x}\sim P_g}\!\left[D(\tilde{x})\right]
      - \mathbb{E}_{x\sim P_r}\!\left[D(x)\right]}_{\text{original critic loss}}
    + \lambda\,
    \underbrace{\mathbb{E}_{\hat{x}\sim P_{\hat{x}}}\!\left[
      \bigl(\lVert \nabla_{\hat{x}} D(\hat{x}) \rVert_2 - 1\bigr)^2
    \right]}_{\text{gradient penalty}},
  \label{eq:methods:healthgan:wgangp}
\end{equation}
The gradient penalty is evaluated at random points lying on the straight
line between a real record $x$ and a generated record $\tilde{x}$. Each
point is calculated as
$\hat{x}=\epsilon x+(1-\epsilon)\tilde{x}$, where
$\epsilon$ is sampled uniformly between $0$ and $1$ for each pair of
records. Values of $\epsilon$ near $1$ place $\hat{x}$ closer to the real
record, whereas values near $0$ place it closer to the generated record.
The penalty encourages the critic's gradient norm at these points to remain
close to $1$.

The first (braced) term is the negated Wasserstein-1 critic objective, where
minimising it maximises the Wasserstein-1 distance estimate between $P_r$
and $P_g$. The second term regularises the critic towards the $1$-Lipschitz
constraint that the Wasserstein formulation requires, penalising departures
of the critic's input-gradient norm from unity. The penalty weight is fixed
at $\lambda = 10$, used throughout the
experiments in the WGAN-GP study~\cite{Gulrajani2017WGANGP}.
The generator is trained to minimise
$-\mathbb{E}_{\tilde{x}\sim P_g}\!\left[D(\tilde{x})\right]$, driving its
samples towards high critic scores. This
gradient-penalty mechanism replaces the weight clipping of the original
Wasserstein GAN~\cite{Arjovsky2017WGAN}. In the local DP-WGAN prototype used in this thesis, 
the gradient-penalty
implementation could not be combined with the required per-example gradient
tracking without changing the original WGAN-GP objective.

\subsubsection{Hyperparameters and Training Configuration}
\label{sec:methods:healthgan:hp}

The training configuration follows the values fixed in the WGAN-GP
implementation. In the implementation, each \emph{epoch} consists of five critic updates
followed by one generator update. Each critic update uses a separate
mini-batch of $11{,}400$ records, so the critic processes $57{,}000$ of the
$57{,}214$ training records during one epoch. An epoch therefore corresponds
approximately to one complete pass through the training data for the critic.
The generator is updated once per epoch using the critic's scores for newly
generated samples. Both networks are optimised
with Adam ($\eta = 10^{-4}$, $\beta_1 = 0.5$, $\beta_2 = 0.9$). The mini-batch
size is derived from the training-set size as
$\lfloor N_{\mathrm{train}}/n_{\mathrm{critic}}\rfloor$, rounded down to the
nearest multiple of $100$, which yields $11{,}400$ for the
$N_{\mathrm{train}} = 57{,}214$ encoded training records. The large batch
follows HealthGAN's design intent of reliably capturing rare clinical
values~\cite{healthgan}. Training runs for $100{,}000$ such iterations, and
the critic loss on a shuffled $11{,}400$-row batch of the held-out test
partition is logged every $1{,}000$ iterations as a training diagnostic.
Table~\ref{tab:healthgan-hp} summarises the configuration.

\begin{table}[h]
\centering
\small
\caption{HealthGAN (WGAN-GP) architecture and training configuration.}
\label{tab:healthgan-hp}
\begin{tabular}{ll}
\toprule
Setting & Value \\
\midrule
Latent dimension & $100$ \\
Encoded feature count $F$ & $45$ \\
Base critic width $B$ & $64$ \\
Generator hidden widths & $90,\ 68$ \\
Critic hidden widths & $64,\ 128,\ 256$ \\
Critic updates per generator update & $5$ \\
Gradient-penalty weight $\lambda$ & $10$ \\
Optimiser & Adam \\
Learning rate $\eta$ & $1\times10^{-4}$ \\
Adam $(\beta_1,\beta_2)$ & $(0.5,\ 0.9)$ \\
Mini-batch size & $11{,}400$ \\
Training iterations & $100{,}000$ \\
Training / test records & $57{,}214 \,/\, 14{,}304$ \\
\bottomrule
\end{tabular}
\end{table}

\subsubsection{Selection of the Synthetic Sample}
\label{sec:methods:healthgan:variants}

Sampling the trained generator is stochastic. A single training run yields
ten interchangeable synthetic files drawn from the same fitted model, of
which three (indexed $0$, $3$, and $6$) were retained for the
comparison below. These files are not identical as each is an
independent draw and therefore exhibits sampling variation in its marginal
and joint statistics. To base the cross-method comparison on a justified
rather than an arbitrary sample, the candidates were evaluated against the
real training partition along two complementary axes: i) marginal fidelity:
quantified by the mean Kolmogorov-Smirnov statistic across columns, and
ii) dependency fidelity: quantified by the mean absolute difference between the
synthetic and real pairwise correlation matrices, with a principal-component
projection used as a visual cross-check of distributional overlap
(Appendix~\ref{app:hg-variants}).

Table~\ref{tab:hg-variants} reports the two fidelity measures for the three
candidates where the spread is negligible. The mean KS statistic ranges over only
$0.0477$-$0.0484$ (a span of $0.0007$) and the mean correlation difference
over $0.0336$-$0.0342$ (a span of $0.0006$). Similarity to the real data is
therefore carried near-identically across the independent draws and no
candidate is meaningfully better than the others. The combined fidelity rank
places draws $0$ and $3$ jointly first (average rank $1.50$). Here, draw $3$ has
the marginally lower KS statistic while draw $0$ has the lowest correlation
difference. Given this tie and the negligible overall spread, draw $0$ was
carried forward to the cross-method comparison as a representative sample.
The comparison indicates that selecting draw $0$ is unlikely to materially
affect the reported marginal and pairwise-correlation fidelity results.

\begin{table}[H]
\centering
\small
\caption{Fidelity of the three retained HealthGAN synthetic draws against
the real training partition (lower is better on both measures). The spread
across draws is negligible. Draw $0$ was selected as representative.}
\label{tab:hg-variants}
\begin{tabular}{lccc}
\toprule
Draw & Mean KS  & Mean $|$corr.\ diff.$|$ & Avg.\ rank \\
\midrule
HealthGAN-0 & 0.0480 & \textbf{0.0336} & 1.50 \\
HealthGAN-3 & \textbf{0.0477} & 0.0340 & 1.50 \\
HealthGAN-6 & 0.0484 & 0.0342 & 3.00 \\
\bottomrule
\end{tabular}
\end{table}

\subsubsection{Rationale for Excluding a DP-Augmented HealthGAN Variant}
\label{sec:methods:healthgan:nodp}

As described in \S\ref{sec:methods:healthgan:loss}, HealthGAN is built on the
WGAN-GP objective, in which a gradient penalty regulates the critic. 
A DP-augmented HealthGAN variant would
therefore need to retain this objective while introducing a formal privacy
mechanism.

A differentially private version of HealthGAN could, in principle, be attempted
by training its critic with DP-SGD. This thesis does not include such a variant
because the local prototype did not preserve the training objective of the
HealthGAN model evaluated in this study.

DP-SGD computes and clips the gradient of each training example to limit
its influence. Gaussian noise is then added to the aggregated clipped
gradients to provide privacy. A privacy accountant then
tracks the total privacy cost across training
iterations~\cite{Abadi2016DPSGD}. WGAN-GP uses a different type of gradient.
Its gradient penalty measures the critic's gradient with respect to
interpolated input samples and encourages the critic to satisfy the Lipschitz
condition required by Wasserstein
training~\cite{Arjovsky2017WGAN,Gulrajani2017WGANGP}. DP-SGD and WGAN-GP
therefore serve different purposes. DP-SGD protects individual training
records, whereas the WGAN-GP penalty regulates the behaviour of the critic.

The local prototype used Opacus to apply DP-SGD to the
critic~\cite{Yousefpour2021Opacus}. However, retaining the gradient 
penalty would require per-example
differentiation through the critic's input-gradient computation. The
selected implementation did not support this operation. 
The prototype therefore removed the gradient penalty and used
the weight-clipping objective of the original
WGAN~\cite{Arjovsky2017WGAN}. As a result, it represented a DP-WGAN rather
than a differentially private version of the HealthGAN WGAN-GP model.

The training requirements also differed substantially. HealthGAN uses large
critic batches so that rare and high-variance clinical values are represented
during training~\cite{healthgan}. DP-SGD requires per-example gradient
calculation~\cite{Abadi2016DPSGD}. Naive implementations of this calculation
can reduce accelerator utilisation and substantially increase training
time~\cite{lee2021scaling}. It also increases memory requirements because a
separate gradient must be retained for each example in a
batch~\cite{Yousefpour2021Opacus}. Opacus provides vectorised
per-sample-gradient computation and mechanisms for separating physical and
logical batch sizes, but these mechanisms still produce a different training
profile from the original HealthGAN
implementation~\cite{Yousefpour2021Opacus}.

For these reasons HealthGAN is retained as the non-private WGAN-GP baseline,
while PATE-GAN is used as the formally private GAN comparator. PATE-GAN was
designed around noisy aggregation of teacher predictions rather than direct
perturbation of critic gradients~\cite{pategan,Papernot2017PATE}. The PATE-GAN
study compares this approach with DPGAN~\cite{Xie2018DPGAN} and reports higher
average predictive utility across the tested privacy
budgets~\cite{pategan}. The comparison between HealthGAN and PATE-GAN should
therefore be interpreted as a comparison between representative non-private
and private GAN approaches. It is not a controlled ablation that isolates the
effect of differential privacy from differences in model architecture and
training procedure.

\subsection{PATE-GAN: Differentially Private GAN}
\label{sec:methods:pategan}

% Code references:
%   - Pategan/pate_gan.py
%   - Pategan/generate_pategan_synthetic.py

PATE-GAN provides the differentially private GAN baseline using a standard
generator, a student discriminator, and an ensemble of teacher 
discriminators.\footnote{The implementation used in this thesis was adapted
from the authors' PATE-GAN GitHub repository
~\cite{jordon2019pateganCode},
\url{https://github.com/vanderschaarlab/mlforhealthlabpub/tree/main/alg/pategan},
and modified for the dataset and experimental pipeline used in this study.}
The teachers are trained on separate, non-overlapping partitions of the real
training data, while the student receives information about the real data only
through the teachers' aggregated predictions~\cite{pategan,Papernot2017PATE}.
Differential privacy is enforced exclusively at the teacher-to-student
interface through Laplace-noised teacher-vote aggregation, so that all subsequent
updates to the student and the generator are post-processing of a
$(\varepsilon, \delta)$-DP signal and therefore inherit the same privacy
guarantee. Unlike the SDV-encoded representation used by HealthGAN
(\S\ref{sec:methods:healthgan:sdv}), PATE-GAN operates directly on the
feature-engineered table. Each of the $D=45$ columns in the 
preprocessed training split is min-max
scaled to $[0,1]$. The generator is trained on this scaled representation.
After generation, the synthetic values are clipped to $[0,1]$ and transformed
back to their original scales. Columns represented as integer-valued variables
in the preprocessed table are then rounded to integer values.

\subsubsection{Teacher-Student Framework}
\label{sec:methods:pategan:teacher}

The training data is partitioned uniformly at random into $k = 10$ disjoint
subsets of equal size. At the beginning of each outer training round, a new logistic-regression
teacher is fitted to each data partition. The teachers are therefore linear
classifiers rather than neural networks. Each teacher is trained on a balanced binary
discrimination task in which its assigned real partition supplies the
positive class and an equal-sized batch sampled from the current generator
supplies the negative class. Because each used real record can affect at 
most one teacher vote, the
sensitivity of the released teacher aggregation is bounded. The student
never observes real data directly. The student is a small
feed-forward network (architecture in \S\ref{sec:methods:pategan:arch})
whose labels are obtained from the teacher ensemble through the noisy
aggregation mechanism of \S\ref{sec:methods:pategan:laplace}, and the
generator is updated against the student through a Wasserstein-style
objective with weight clipping~\cite{Arjovsky2017WGAN}.

\subsubsection{PATE Voting with the Laplace Mechanism}
\label{sec:methods:pategan:laplace}

For an input feature vector $x$, let $T_i$ denote the $i$-th teacher,
with $i=1,\dots,k$, and let $m$ denote the number of possible classes.
The PATE-GAN paper defines the teacher vote count for class $j$ as
\begin{equation}
  n_j(x) \;=\; \bigl|\{T_i : T_i(x) = j\}\bigr|,
  \qquad j = 1,\dots,m,
\end{equation}
where $n_j(x)$ is the number of teachers assigning class $j$ to
$x$~\cite{pategan}. The PATE mechanism adds independent Laplace noise to each class count and
returns the class with the largest noisy count,
\begin{equation}
  \mathrm{PATE}_{\lambda}(x)
  \;=\;
  \operatorname*{arg\,max}_{j \in \{1,\dots,m\}}
  \bigl(n_j(x) + Y_j\bigr),
  \qquad
  Y_1,\dots,Y_m \overset{\mathrm{i.i.d.}}{\sim} \mathrm{Lap}(\lambda).
  \label{eq:pategan-aggregation}
\end{equation}
This noisy-argmax mechanism is analysed in the PATE framework and adopted
by PATE-GAN~\cite{Papernot2017PATE,pategan}. In this thesis, the teachers
assign binary real or generated labels, so $m=2$, and the Laplace noise
scale is set to $\lambda=1.0$.

\subsubsection{Privacy Accounting via the Moments Accountant}
\label{sec:methods:pategan:accountant}

Each release of a noisy teacher vote contributes to the total privacy cost.
PATE-GAN tracks this cumulative cost using the moments accountant. 
Queries for which the teachers strongly agree incur a lower
privacy cost than queries with closely divided votes~\cite{pategan}.

After each outer training round, the accumulated privacy cost is converted
into an $(\varepsilon,\delta)$ guarantee:
\begin{equation}
  \hat{\varepsilon}
  =
  \min_{l \in \{1,\dots,L\}}
  \frac{\alpha(l)+\log(1/\delta)}{l},
  \label{eq:pategan-epshat}
\end{equation}
where $\alpha(l)$ is the accumulated privacy-loss moment of order $l$ and
$L$ is the number of tracked orders~\cite{pategan}. Training continues while
$\hat{\varepsilon}$ remains below the target $\varepsilon$ and stops when
the target is reached or exceeded. Thus, the privacy budget determines the
number of training rounds instead of a fixed number of epochs. The final
synthetic dataset is then generated from the trained model.

\subsubsection{Architecture}
\label{sec:methods:pategan:arch}

The generator and student are fully connected neural networks. Their
dimensions are based on the $D=45$ features in the preprocessed table.

The generator receives a random vector
$z \sim \mathcal{U}([-1,1]^D)$ and transforms it through three dense layers:
\begin{equation}
  D \;\to\; 4D~(\mathrm{tanh}) \;\to\; 4D~(\mathrm{tanh})
  \;\to\; D~(\mathrm{sigmoid}).
\end{equation}
The hidden layers therefore contain $4D=180$ units. The sigmoid output
restricts each generated feature to $[0,1]$, matching the min-max scaled
training data.

The student assigns each record an unrestricted scalar score using:
\begin{equation}
  D \;\to\; D~(\mathrm{ReLU}) \;\to\; 1.
\end{equation}
The output layer has no activation function, as required by the Wasserstein
formulation~\cite{Arjovsky2017WGAN}. The generator and student parameters are
initialised using a layer-width-scaled normal distribution. The teachers are
logistic-regression models with default regularisation, as described in
\S\ref{sec:methods:pategan:teacher}.

\subsubsection{Privacy Budget and Hyperparameters}
\label{sec:methods:pategan:hp}

The student and generator are both optimised with 
Root Mean Square Propagation (RMSProp) at learning
rate $\eta = 10^{-4}$. The student is trained under the Wasserstein
formulation with hard weight clipping. After every gradient step its
parameters are projected back into the box $[-0.01, 0.01]$, which
enforces the $1$-Lipschitz constraint required of the
critic~\cite{Arjovsky2017WGAN}. The target privacy budget is
$(\varepsilon, \delta) = (5.0,\,10^{-5})$, matched to the budget used
for the Pythia-1B-DP run (\S\ref{sec:methods:pythia-dp:epsilon}), so that
the GAN-side and LLM-side DP comparisons in
\S\ref{sec:results:rq3:gan-toggle} and \S\ref{sec:results:rq3:llm-toggle}
are reported at a common $\varepsilon$. Training runs until the
moments-accountant estimate $\hat{\varepsilon}$ from
Equation~(\ref{eq:pategan-epshat}) reaches this target. After
termination, one synthetic batch of size $N_{\mathrm{train}} = 57{,}214$
is drawn from the generator and inverted from the min-max
representation to the original feature space. Table~\ref{tab:pategan-hp}
summarises the configuration.

\begin{table}[h]
\centering
\small
\caption{PATE-GAN architecture, training, and privacy configuration}
\label{tab:pategan-hp}
\begin{tabular}{ll}
\toprule
Setting & Value \\
\midrule
Feature count $D$ & $45$ \\
Latent distribution & $\mathcal{U}([-1,1]^{D})$ \\
Latent dimension $z_{\mathrm{dim}}$ & $45$ \\
Generator hidden widths & $180,\ 180$ \\
Student hidden width & $45$ \\
Teacher classifier & Logistic regression (linear) \\
Number of teachers $k$ & $10$ \\
Records per teacher partition & $5{,}721$ \\
Student updates per outer round $n_s$ & $1$ \\
Mini-batch size & $64$ \\
Laplace noise scale $\lambda$ & $1.0$ \\
Moments-accountant orders $L$ & $20$ \\
Optimiser & RMSProp \\
Learning rate $\eta$ & $1\times 10^{-4}$ \\
Student weight-clipping range & $[-0.01,\ 0.01]$ \\
Target privacy budget $(\varepsilon,\delta)$ & $(5.0,\ 10^{-5})$ \\
\bottomrule
\end{tabular}
\end{table}

\subsection{Pythia-70M: Non-DP LLM (Small-Scale Baseline)}
\label{sec:methods:pythia}

% Code references:
%   - pythia/pythia_tabular.py
%   - pythia/generate_pythia_synthetic.py

Pythia-70M realises the non-private LLM baseline at the small end of the
scale axis. The pipeline serialises each tabular row as a single text
sequence with a class-label prefix. It then fine-tunes the base model 
on those sequences using low-rank
adaptation. To generate synthetic data, the trained model is prompted with the same
class prefix used during training. The generated text is then converted
into typed values and validated against the source schema. The following
subsubsections describe the base model, row-to-text serialisation, LoRA
fine-tuning, class-conditioned generation, and schema-aware
post-processing.

\subsubsection{Base Model}
\label{sec:methods:pythia:base}

The base model is Pythia-70M, a decoder-only transformer with $70$M
parameters drawn from the open-weights Pythia
suite~\cite{Biderman2023Pythia}. The choice of the smallest Pythia
variant is dictated by training-cost considerations. It provides a compute-tractable 
LLM baseline as the smallest official Pythia scale while
retaining the same model family and tokenisation scheme as the larger
Pythia variants. Pythia-based Tabula
models in some settings outperforms the CTGAN
reference baseline on TSTR utility when
applied to electronic-health-record tables~\cite{Miletic2024}. 
This motivates including the smallest variant as a
non-private baseline rather than excluding it on capacity grounds.
Tokenisation uses the model's native byte-pair-encoded vocabulary. The
end-of-sequence token is reused as the padding token when the
pretrained tokeniser lacks an explicit one, and the padding side is set
to the left to preserve the structure of generation prompts.

\subsubsection{Tabular-to-Text Serialisation}
\label{sec:methods:pythia:serialisation}

Each row of the preprocessed training table is serialised into a single
text sequence of the form
\begin{equation*}
  \mathrm{Class}\_\langle L\rangle ~|~ c_1 = v_1 ~|~ c_2 = v_2 ~|~ \dots ~|~ c_K = v_K,
\end{equation*}
where $\langle L \rangle$ is the integer class label drawn from the
binary readmission target column, $c_i$ is the $i$-th column name in
the original dataframe order, and $v_i$ is the textual representation
of that column's value in the row. Scalar values are formatted
deterministically where integers retain their textual decimal form,
floating-point values are formatted with $15$ significant digits,
booleans are mapped to $\{0, 1\}$, and missing values are emitted as
the literal string ``NA''. The class prefix is fixed at the head of
the sequence so that conditional generation can later be triggered by
pinning that prefix to the desired label
(\S\ref{sec:methods:pythia:generation}).

\subsubsection{LoRA Fine-tuning}
\label{sec:methods:pythia:lora}

The base Pythia weights are kept frozen and adapted through low-rank
adaptation (LoRA)~\cite{Hu2022LoRA}. The LoRA configuration uses rank
$r = 8$, scaling factor $\alpha = 16$, dropout $0.05$, no bias
adaptation, and the causal-language-modelling task type. So only
the injected low-rank matrices receive gradient updates. The corpus is
tokenised without padding and dynamically padded per mini-batch to a
multiple of eight by a standard causal-LM data collator. 
Training runs for $10$ epochs at per-device batch size $8$,
learning rate $\eta = 2 \times 10^{-5}$, and maximum sequence length
$512$, with half-precision arithmetic on CUDA-capable devices. The 
mean training loss is recorded for each epoch. Whenever an epoch 
achieves a new minimum loss, the trainable parameters are snapshotted. 
At the end of training, the snapshot with the lowest recorded training
loss is restored for generation. The
fine-tuned model is moved into evaluation mode before generation.

\subsubsection{Class-Conditioned Generation}
\label{sec:methods:pythia:generation}

Generation is performed class-by-class so that the synthetic table
preserves the per-class row counts of the source split. For each
target class $\ell \in \{0, 1\}$, the model is prompted with a class
prefix of the form $\mathrm{Class}\ \ell ~|~ c_1 =$. Here, $c_1$ is
the name of the first non-target column under the dataframe's column
order, and is left to autoregressively complete the remainder of the
row. Generation uses nucleus sampling with a temperature of $0.8$ and a
top-$p$ threshold of $0.95$. The model generates between $64$ and $512$
new tokens to reduce row truncation, with the key-value cache enabled. 
The required class label is enforced after parsing
rather than inferred from the generation, removing any dependence of
the assigned label on the model's own output
(\S\ref{sec:methods:pythia:parsing}).

\subsubsection{Schema-Aware Postprocessing}
\label{sec:methods:pythia:schema}
\label{sec:methods:pythia:parsing}

Before training, a schema is derived from the source data frame and
stored for processing the generated outputs. For each column, the schema
records the original pandas data type, whether the values are numerical,
and, for numerical columns, the observed minimum and maximum. It also records 
whether all observed values are integer-valued, using this as the 
integer-coded flag. When a numeric column has at most 25 distinct values, 
the schema stores its full sorted discrete support, which includes 
ordinal-coded columns such as ten-year age bands and diagnosis 
groupings as well as binary medication 
flags. For non-numeric columns, it stores the full list of distinct string 
categories and the column mode used as a schema-safe fallback when a generated 
field cannot be parsed.

Generated text is then parsed into key-value fields and coerced
against this schema. Numeric fields
are clipped to their observed ranges, snapped to the recorded discrete
support where applicable, and rounded when the source column is
integer-coded. Non-numeric fields are checked against the recorded
category set, with invalid or missing values replaced by schema-safe
fallbacks. The target column is not trusted from the model output. It
is overwritten with the externally requested class label so that the
generated class matches the conditioning prompt.

Generation is retried until the requested number of rows for each
class is accepted or the fixed retry budget is exhausted. If too few
valid rows are obtained, the implementation fills the shortfall from
already accepted rows for that class, with a schema-valid random
fallback available only as a last resort. The final data frame is 
shuffled, restored to the original column order,
and converted to the original data types. It is then validated for row
count, missing values, target values, and numerical bounds. Finally, the
data and metadata describing the acceptance and retry statistics are
saved.

\subsection{Pythia-1B: Non-DP LLM (Scale Anchor)}
\label{sec:methods:pythia-1b}

% Code references:
%   - pythia/pythia_tabular.py (run at -model_name EleutherAI/pythia-1b)
%   - pythia/generate_pythia_synthetic.py

A second non-private LLM variant is trained at a larger scale to serve as
the non-DP anchor of the LLM-side DP toggle
(\S\ref{sec:methods:pythia-dp}). Pythia-1B uses the same row-to-text
serialisation (\S\ref{sec:methods:pythia:serialisation}), class-conditioned
generation procedure (\S\ref{sec:methods:pythia:generation}), and
schema-aware postprocessing (\S\ref{sec:methods:pythia:parsing}) as the 70M
variant. It also uses
the same non-private LoRA adapter structure, while run-level optimisation
settings are adjusted for the larger backbone. The two non-private variants 
use the same tabular representation and
postprocessing pipeline, while differing in backbone scale and the
run-level optimisation settings required for each model. There was no change made
in tabular representation or postprocessing logic~\cite{Biderman2023Pythia}.
Including Pythia-1B serves two purposes.
First, it isolates the cost of DP-SGD in the RQ3 within-family comparison
(Pythia-1B versus Pythia-1B-DP) from the confounding effect of a simultaneous
model-scale change. Second, it positions an LLM-side scale axis in the RQ2
landscape view, enabling a qualitative read of how the LLM family's
privacy-utility profile shifts with model size under a shared tabular
representation and postprocessing pipeline.

\subsection{Pythia-1B-DP: Differentially Private LLM}
\label{sec:methods:pythia-dp}

% Code references:
%   - pythia/pythia_tabular_dp.py
%   - pythia/generate_pythia_synthetic_dp.py

This subsection describes the differentially private LLM used in the RQ3
privacy comparison. It follows the same tabular serialisation,
class-conditioned generation, and schema-aware post-processing procedures
as the non-private Pythia pipeline. However, the LoRA adapters are trained
using DP-SGD through Opacus.

\subsubsection{Rationale for Operationalising DP-SGD at the 1B Scale}
\label{sec:methods:pythia-dp:why-1b}

An initial attempt to apply DP-SGD to Pythia-70M at $\varepsilon = 5.0$
failed to train usefully. The recorded training loss worsened across the
five epochs (from $3.38$ to $3.74$), and the run was not carried forward
as an evaluable synthetic-table generator.
Scaling the LLM-side DP variant to Pythia-1B under the same
target privacy budget and otherwise comparable tabular-generation pipeline
restored stable convergence under the same
$(\varepsilon = 5.0, \delta = 10^{-5})$ budget. This outcome is consistent
with the observation that larger
LLMs tolerate DP noise more gracefully, and motivates the asymmetric LLM
lineup~\cite{Li2022LargeDP}. Although it's important to note that the batch size, learning rate 
and noise multiplier was adjusted for a better outcome. The LLM-side DP toggle is operationalised at the 1B scale, while
the 70M baseline is retained only as a non-DP small-scale reference in
the RQ2 landscape view. The empirical failure itself is reported as a
finding in \S\ref{sec:results:rq3:pythia-70m-failure}.

\subsubsection{Private LoRA Training with DP-SGD}
\label{sec:methods:pythia-dp:training}

Differential privacy is implemented using differentially private
stochastic gradient descent (DP-SGD)~\cite{Abadi2016DPSGD}. The Opacus
library provides per-sample gradient computation, noise calibration, and
privacy accounting~\cite{Yousefpour2021Opacus}.
The pipeline uses parameter-efficient fine-tuning, keeping the base Pythia
weights fixed and applying differential privacy only to the LoRA adapter
parameters. By the post-processing property of differential privacy, the
resulting generator satisfies $(\varepsilon,\delta)$-DP with respect to the
private training data~\cite{Yu2022DPFT}. The AdamW 
optimiser is restricted to the trainable adapter parameters. The
optimiser, model, and data loader are then passed to the Opacus privacy engine which 
computes and clips per-sample gradients before adding Gaussian noise.
Its built-in accountant calibrates the noise multiplier using the target
$(\varepsilon,\delta)$ budget and the planned number of training epochs.

Following the clipping and noise-addition mechanism defined in
Eqs.~\ref{eq:bg:dp:clip}-\ref{eq:bg:dp:update}, Opacus computes a gradient
for each training example and clips its $\ell_2$ norm to $C=1.0$. The clipped
gradients are then summed, and Gaussian noise with standard deviation
$\sigma C$ is added. This noisy sum is averaged over the logical batch and
used by AdamW to update the model.

Opacus automatically selects the noise multiplier $\sigma$ according to the
target $(\varepsilon,\delta)$ budget and the planned number of training
epochs. In the DP configuration, LoRA dropout is set to zero, compared with
$0.05$ in the non-DP configuration
(\S\ref{sec:methods:pythia:lora}). All other LoRA settings, including the rank,
scaling factor, task type, and bias configuration, remain unchanged.

\subsubsection{Privacy Budget and Pipeline Differences}
\label{sec:methods:pythia-dp:epsilon}%
\label{sec:methods:pythia-dp:diff}%

The privacy accountant uses the \emph{logical} batch size rather than the
physical batch processed by the device. The logical batch size is
\begin{equation*}
  B_{\mathrm{eff}}
  = B_{\mathrm{phys}} \times G
  = 32 \times 16
  = 512,
\end{equation*}
where $B_{\mathrm{phys}}=32$ is the physical batch size and $G=16$ is
the gradient-accumulation factor. Opacus divides each logical batch into
smaller physical batches of 32 records. These smaller batches are processed
separately to limit the memory required for per-sample gradients. For privacy
accounting, Opacus uses the logical-batch sampling rate
$q=B_{\mathrm{eff}}/N_{\mathrm{train}}$.

All DP-enabled Pythia runs use a target privacy budget of
$(\varepsilon,\delta)=(5.0,\,10^{-5})$. The value of $\delta$ is set well
below $1/N_{\mathrm{train}}$, following standard practice
\cite{Abadi2016DPSGD}. Opacus determines the Gaussian noise multiplier
$\sigma$ from the privacy budget, the number of training epochs, and the
sampling rate $q$.

After training, the privacy accountant calculates the realised
$\varepsilon$ at the selected value of $\delta$. The implementation records
this value together with the noise multiplier, sampling rate, loss for each
epoch, and best-performing epoch. These values are stored in a metadata file
alongside the synthetic dataset and reported in
Chapter~\ref{ch:results}.

PATE-GAN uses the same target privacy budget
(\S\ref{sec:methods:pategan:hp}). This common budget supports the comparison
of the GAN and LLM privacy settings in
\S\ref{sec:results:rq3:gan-toggle} and
\S\ref{sec:results:rq3:llm-toggle}.

Except for the DP-SGD training procedure, the DP variant uses the same
pipeline as the non-DP Pythia model described in
\S\ref{sec:methods:pythia}. The serialisation, schema derivation, sampling,
parsing, and validation procedures remain unchanged. The larger logical
batch reduces variation in the averaged noisy gradient, while Opacus
calibrates the noise multiplier to satisfy the target privacy budget.
Table~\ref{tab:pythia-dp-diff} summarises the remaining differences between
the two pipelines.

\begin{table}[h]
\centering
\small
\caption{Differences of the DP pipeline from the non-DP Pythia pipeline
(\S\ref{sec:methods:pythia}). All other logic is shared.}
\label{tab:pythia-dp-diff}
\begin{tabular}{lll}
\toprule
Aspect & Non-DP path & DP path \\
\midrule
Training loop & Trainer driver & Manual epoch loop over Opacus-wrapped loader \\
LoRA dropout & $0.05$ & $0.0$ (deterministic forward for per-sample grads) \\
Batch sizing & $32$, no DP accumulation & $B_{\mathrm{phys}}{=}32 \times G{=}16 = 512$ logical \\
Memory & Standard & Per-sample grads $\times B_{\mathrm{phys}}$; micro-batched loader \\
Inference & Direct generation & Unwrap per-sample-gradient module before sampling \\
\bottomrule
\end{tabular}
\end{table}

\subsection{Implementation Summary}
\label{sec:methods:summary}

The five generators introduced in \S\S\ref{sec:methods:healthgan}-%
\ref{sec:methods:pythia-dp} share a common preprocessing front end
(\S\ref{sec:method:data}) and a common evaluation back end
(\S\ref{sec:method:evaluation}), but differ along two main
implementation axes: i) the generative paradigm and ii) the privacy mechanism.
Table~\ref{tab:methods-implementation-summary} consolidates the
configuration of every generator along these axes for ease of
cross-reference.

\begin{table}[h]
\centering
\footnotesize
\caption{Comparative summary of the five synthesisers. All generators
are trained on the same $57{,}214$-record training partition of the
Diabetes 130-US Hospitals benchmark and supply synthetic outputs whose
sizes are matched to the partitions used in the downstream evaluation
pipeline of Chapter~\ref{ch:results}. Privacy budgets are reported as
target $(\varepsilon, \delta)$ values. The realised $\varepsilon$ for
Pythia-1B-DP is reported in \S\ref{sec:results:rq3:llm-toggle}.}
\label{tab:methods-implementation-summary}
\begin{tabular}{llll}
\toprule
Generator & Family & Privacy mechanism & $(\varepsilon, \delta)$ \\
\midrule
HealthGAN
  & GAN (WGAN-GP)
  & None
  & n/a \\
PATE-GAN
  & GAN (PATE + WGAN loss)
  & PATE Laplace aggregation
  & $(5.0,\,10^{-5})$ \\
Pythia-70M-LoRA
  & LLM (decoder-only)
  & None
  & n/a \\
Pythia-1B-LoRA
  & LLM (decoder-only)
  & None
  & n/a \\
Pythia-1B-DP-LoRA
  & LLM (decoder-only)
  & DP-SGD on LoRA adapter
  & $(5.0,\,10^{-5})$ \\
\bottomrule
\end{tabular}
\end{table}

Two aspects of the experimental design are important for the comparison.
First, the models apply differential privacy at different stages of training.
PATE-GAN adds calibrated Laplace noise when teacher predictions are combined
to label data for the student
\cite{pategan,Papernot2017PATE}. In contrast, Pythia-1B-DP applies DP-SGD
by clipping the per-sample gradients of the LoRA adapter and adding Gaussian
noise~\cite{Abadi2016DPSGD,Yu2022DPFT}.

Both methods use the target privacy budget
$(\varepsilon,\delta)=(5.0,\,10^{-5})$. This common budget provides a
consistent reference for the GAN and LLM comparisons reported in
\S\ref{sec:results:rq3:gan-toggle} and
\S\ref{sec:results:rq3:llm-toggle}.

Second, all Pythia variants use the same row-to-text serialisation and
schema-aware parsing procedures. The non-DP variants also share the same
LoRA adapter structure. The DP variant retains the LoRA rank, scaling
factor, task type, and bias configuration. However, it sets LoRA dropout
to zero and introduces DP-SGD gradient clipping and noise addition
(\S\ref{sec:methods:pythia-dp:diff}).